\documentclass[10pt]{article}
\usepackage[T1]{fontenc}
\usepackage[utf8]{inputenc}
\usepackage[french]{babel}
\usepackage{amsmath,amssymb,mathrsfs,mathpazo}
\usepackage{geometry}
\geometry{a4paper,margin=2cm}
\usepackage{graphicx,array,tabularx,tabularray,colortbl,multirow,multicol,booktabs,enumitem,pifont,ulem,fancybox,diagbox,cancel}
\usepackage{pstricks,pstricks-add,pst-plot,pst-tree,pst-node,pst-text,pst-func,pst-math,pst-bspline,pst-3dplot,pst-eucl,pst-solides3d}
\let\pscalculate\undefined
\usepackage{pst-bezier}
\usepackage{tikz,pgfplots,tkz-tab}
\usetikzlibrary{arrows,arrows.meta,patterns,calc,positioning,trees,decorations.markings,shapes.geometric,shapes.misc,shapes.arrows,intersections,decorations.pathmorphing,decorations.pathreplacing,backgrounds,shadows,babel}
\pgfplotsset{compat=1.18}
\usepgfplotslibrary{fillbetween}
\usepackage[np]{numprint}
\usepackage{listings,makecell,pas-tableur,twoopt}
\let\barre\undefined
\usepackage[pstricks]{tablvar}
\definecolor{darkgreen}{RGB}{0,100,0}
\providecommand{\red}{\color{red}}
\providecommand{\blue}{\color{blue}}
\providecommand{\green}{\color{green}}
\providecommand{\black}{\color{black}}
\providecommand{\white}{\color{white}}
\providecommand{\gray}{\color{gray}}
\providecommand{\cyan}{\color{cyan}}
\providecommand{\magenta}{\color{magenta}}
\providecommand{\yellow}{\color{yellow}}
\providecommand{\violet}{\color{violet}}
\providecommand{\orange}{\color{orange}}
\providecommand{\euro}{\texteuro}
\providecommand{\EUR}{\texteuro}
\providecommand{\ieme}{\textsuperscript{e}}
\providecommand{\degres}{\ensuremath{^\circ}}
\providecommand{\degree}{\ensuremath{^\circ}}
\providecommand{\vv}[1]{\overrightarrow{#1}}
\providecommand{\Coord}[1]{#1}
\providecommand{\decosix}{}
\providecommand{\starredbullet}{\ensuremath{\bullet}}
\pagestyle{empty}
\setlength{\parindent}{0pt}
\usepackage{adjustbox}\DeclareUnicodeCharacter{2236}{:}\DeclareUnicodeCharacter{2190}{\ensuremath{\leftarrow}}
\providecommand{\R}{\mathbb{R}}
\providecommand{\N}{\mathbb{N}}
\providecommand{\Z}{\mathbb{Z}}
\providecommand{\Q}{\mathbb{Q}}
\providecommand{\D}{\mathbb{D}}
\providecommand{\C}{\mathbb{C}}
\providecommand{\e}{{\mathrm{e}}}
\providecommand{\vect}[1]{\overrightarrow{#1}}
\providecommand{\vectt}[1]{\overrightarrow{\mathrm{#1}}}
\providecommand{\vvt}[1]{\overrightarrow{\mathrm{#1}}}
\providecommand{\np}[2][]{{#2~#1}}
\providecommand{\vv}[1]{{\overrightarrow{#1}}}
\renewcommand{\d}{{\mathrm{d}}}
\providecommand{\degres}{{^\circ}}
\providecommand{\degre}{{^\circ}}
\providecommand{\blue}{\color{blue}}
\providecommand{\red}{\color{red}}
\providecommand{\enskip}{\quad}
\providecommand{\llbracket}{\lbrack\!\lbrack}
\providecommand{\rrbracket}{\rbrack\!\rbrack}
\providecommand{\euro}{{\text{€}}}
\providecommand{\no}{{\text{n°}}}
\providecommand{\interval}[2]{{[#1;#2]}}
\providecommand{\intervFO}[2]{{[#1;#2[}}
\providecommand{\dotfill}{\ldots}

\providecommand{\ds}{\displaystyle}
\providecommand{\textding}[1]{#1}
\providecommand{\Oij}{\(\mathrm{O};\vec{\imath},\vec{\jmath}\)}
\providecommand{\Oijk}{\(\mathrm{O};\vec{\imath},\vec{\jmath},\vec{k}\)}
\providecommand{\Ouv}{\(\mathrm{O};\vec{u},\vec{v}\)}
\providecommand{\intervFF}[2]{\ensuremath{[#1;#2]}}
\renewcommand{\euro}{\ifmmode\text{\texteuro}\else\texteuro\fi}
\DeclareUnicodeCharacter{00B0}{\ifmmode^{\circ}\else\textdegree\fi}
\DeclareUnicodeCharacter{20AC}{\ifmmode\text{\texteuro}\else\texteuro\fi}
% Énoncés : Éducation nationale. Transcription et corrigés : APMEP, Denis Vergès et contributeurs. Découpage et classement : CoopMaths, Sébastien Lozano, Éric Elter et contributeurs. Adaptation de la présentation : Maths Explorer.
% CC BY-SA : https://coopmaths.fr/www/about/
% https://coopmaths.fr/alea/static/e3c/2020/tex/e3c_2020_00_sujet9_1.tex
\begin{document}
\noindent\textit{APMEP / CoopMaths. CC BY-SA. Présentation adaptée par Maths Explorer.}\par\medskip
\medskip

Ce QCM comprend 5 questions indépendantes. Pour chacune d'elles, une seule des affirmations proposées est exacte.

Indiquer pour chaque question sur la copie la lettre correspondant à la réponse choisie. Aucune justification n'est demandée.

Chaque réponse correcte rapporte $1$ point. Une réponse incorrecte ou une absence de réponse n'apporte ni ne retire de point.

\medskip

\textbf{Question 1}

On choisit au hasard un individu parmi les passagers en transit dans un aéroport. On a représenté ci-dessous un arbre de probabilités lié à certains évènements dont certains éléments ont été effacés.

On considère les évènements suivants:
\begin{itemize}
\item $A$ : \og le passager parle anglais \fg
\item $B$ : \og le passager ne parle pas anglais \fg
\item $E$ : \og le passager est un membre de l'Union Européenne \fg
\end{itemize}

\begin{center}
\pstree[treemode=R,nodesepA=0pt,nodesepB=3pt]{\TR{}}
{\pstree{\TR{$A$~~}\naput{0,6}}
	{\TR{$E$}
	\TR{$\overline{E}$}\tbput{0,5}
	}
\pstree{\TR{$B$~~}}
	{\TR{$E$}\taput{0,3}
	\TR{$\overline{E}$}
	}
}
\end{center}
\begin{center}
\begin{tabularx}{\linewidth}{|*{4}{X|}}\hline
\textbf{a.~~}$P_B(E) = 0,12$&\textbf{b.~~}$p(E) = 0,42$&\textbf{c.~~}\small La probabilité que le passager choisi soit européen et ne parle pas anglais est 0,3&\textbf{d.~~}$P(A \cup B) = 1,1$\\ \hline
\end{tabularx}
\end{center}

\medskip

\textbf{Question 2}

Le plan est muni d'un repère orthonormé. Soit $D$ la droite d'équation $3x + y - 2 = 0$.

\begin{center}
\begin{tabularx}{\linewidth}{|*{4}{X|}}\hline
\textbf{a.~~}Le point de coordonnées $(6~;~-15)$ appartient à $D$&\textbf{b.~~}$D$ est perpendiculaire à la droite d'équation $12x + 4y = 0$&\textbf{c.~~}Le vecteur de coordonnées
(1~;~3) est un vecteur directeur de $D$.&\textbf{d.~~}\small Le vecteur de coordonnées
(3~;~1) est un vecteur directeur des droites perpendiculaires à $D$.\\ \hline
\end{tabularx}
\end{center}

\medskip

\textbf{Question 3} On considère dans l'ensemble des réels l'équation trigonométrique $\sin x = 1$.

\begin{center}
\begin{tabularx}{\linewidth}{|*{4}{X|}}\hline
\textbf{a.~~}Cette équation admet une unique solution dans l'ensemble des réels&\textbf{b.~~}Cette équation admet une infinité de solutions dans l'ensemble des réels&\textbf{c.~~}$2\pi$ est une solution de cette équation&\textbf{d.~~}$- \dfrac{57\pi\rule{0pt}{11pt}}{2\rule[-3pt]{0pt}{0pt}}$ est une solution de cette équation\\ \hline
\end{tabularx}
\end{center}

\medskip

\textbf{Question 4}

Soit $f$ la fonction définie sur l'ensemble des nombres réels par $f(x) = \dfrac{2x}{x^2 + 1}$
et $\mathcal{C}$ sa courbe représentative dans un repère du plan.

\begin{center}
\begin{tabularx}{\linewidth}{|*{4}{X|}}\hline
\textbf{a.~~} La courbe $\mathcal{C}$ n'admet pas de tangente au point d'abscisse $0$&\textbf{b.~~}La tangente à $\mathcal{C}$ au point d'abscisse $0$ pour équation $y = 2x$&\textbf{c.~~} La tangente à $\mathcal{C}$ au point d'abscisse $0$ a pour coefficient directeur 1&\textbf{d.~~}La tangente à $\mathcal{C}$ au point d'abscisse $0$ est parallèle à l'axe des abscisses\\ \hline
\end{tabularx}
\end{center}

\medskip

\textbf{Question 5}

Soit la fonction $f$ définie sur l'intervalle $]-2~;~ +\infty[$ par:

\[f(x) = \dfrac{x - 3}{x+2}\]

$f$ est dérivable sur l'intervalle $]-2~;~ +\infty[$ et pour tout réel $x$ de $]-2~;~ +\infty[$, on a :

\begin{center}
\begin{tabularx}{\linewidth}{|*{4}{X|}}\hline
\textbf{a.~~}$f'(x) = 1$&\textbf{b.~~}$f'(x) = \dfrac{2x - 1}{(x + 2)^2\rule[-5pt]{0pt}{0pt}}$&\textbf{c.~~}$f'(x) = \dfrac{5}{(x + 2)^2}$&\textbf{d.~~}$f'(x) = 2x - 1$\rule[-3mm]{0mm}{8mm}\\ \hline
\end{tabularx}
\end{center}

\bigskip

\end{document}
